\subsubsection{Konzistence}

Požadavky na konzistenci se zaměřují na to, aby uživatelské rozhraní bylo jednotné a předvídatelné jako jiné systémy, které již uživatelé používají \parencite{rejnkovaLokalizacePrizpusobeniMetodiky}. Tyto požadavky jsou popsány v tabulce \ref{tab:KON-pozadavky}.

\begin{OstatniPozadavkyTabulka}{KON}{Požadavky na konzistenci}
  \OstatniPozadavekRow{Jednotná terminologie napříč rozhraním}{Termíny slovníku pojmů (budova, místnost, senzor, pozorování, subjekt zájmu, veličina, jednotka) musí být v uživatelském rozhraní, dokumentaci, chybových zprávách i v API odpovědích používány konzistentně. Pro každý pojem existuje jedna kanonická česká a jedna kanonická anglická forma; synonyma se nepoužívají.}{2}
  \OstatniPozadavekRow{Jednotné chování formulářů}{Všechny formuláře aplikace (registrace uživatele, registrace budovy, místnosti, senzoru, úprava parametrů) musí mít konzistentní chování: validace probíhá inkrementálně při ztrátě fokusu pole, chybové hlášky se zobrazují u příslušného pole, povinná pole jsou jednotně vizuálně označena, akční tlačítko (Uložit / Vytvořit) je vždy umístěné na stejné pozici.}{2}
\end{OstatniPozadavkyTabulka}
